\chapter{Конструкторская часть}

\section{Диаграмма базы данных}

На основе анализа предметной области была спроектирована реляционная схема базы данных. Разработанная структура обеспечивает хранение информации о пользователях, игровых проектах, модификациях, версиях и файлах, а также поддерживает систему категорий, комментариев и пользовательских коллекций. 

Диаграмма базы данных представлена на рисунке~\ref{img:dbml}.

\begin{figure}[H]
	\centering
	\includegraphics[width=\textwidth]{dbml}
	\caption{Диаграмма базы данных}
	\label{img:dbml}
\end{figure}

\section{Ролевая модель на уровне базы данных}

Для разграничения прав доступа в системе выделены три роли. Матрица доступа к таблицам представлена в таблице~\ref{tab:roles}.
В таблице были использованы следующие обозначения:
\begin{itemize}
	\item C (Create) --- создание;
	\item R (Read) --- чтение;
	\item U (Update) --- обновление;
	\item D (Delete) --- удаление.
\end{itemize}
\renewcommand{\arraystretch}{1.1}

\begin{table}[H]
	\centering
	\caption{Ролевая модель доступа к объектам БД}
	\label{tab:roles}
	\begin{tabularx}{\textwidth}{|X|X|X|X|}
		\hline
		\textbf{Таблица} & \textbf{Гость} & \textbf{Пользователь} & \textbf{Администратор} \\
		\hline
		users & R & RU & CRUD \\
		\hline
		games & R & R & CRUD \\
		\hline
		mods & R & CRUD & CRUD \\
		\hline
		mod\_versions & R & CRUD & CRUD \\
		\hline
		mod\_files & R & CRD & CRUD \\
		\hline
		categories & R & R & CRUD \\
		\hline
		mods\_categories & R & CRD & CRUD \\
		\hline
		comments & R & CRD & CRUD \\
		\hline
		collections & R & CRUD & CRUD \\
		\hline
		collections\_mods & R & CRUD & CRUD \\
		\hline
	\end{tabularx}
\end{table}

\section{Анализ функциональных зависимостей и проверка третьей нормальной формы}

Для исключения избыточности и аномалий обновления проведена нормализация схемы. Ниже приведены описания атрибутов отношений, выделены потенциальные ключи и проведен анализ функциональных зависимостей (ФЗ) на соответствие условиям третьей нормальной формы (3НФ).

\subsection*{Отношение users}

\begin{table}[H]
	\centering
	\caption{Атрибуты отношения users}
	\label{tab:users}
	\begin{tabularx}{\textwidth}{|l|l|X|}
		\hline
		\textbf{Атрибут} & \textbf{Тип} & \textbf{Назначение} \\
		\hline
		username (PK) & Строка & Уникальный логин пользователя \\
		\hline
		name & Строка & Отображаемое имя \\
		\hline
		email & Строка & Электронная почта \\
		\hline
		password & Строка & Хэш пароля \\
		\hline
		register\_date & Дата и время & Дата регистрации \\
		\hline
		is\_admin & Логический & Признак администратора \\
		\hline
	\end{tabularx}
\end{table}

\textit{Потенциальные ключи:} $K_1 = \{username\}$, $K_2 = \{email\}$, $K_3 = \{name\}$. \\
\textit{ФЗ:} $F = \left\{ \begin{array}{l}
	username \rightarrow \{name, email, password, register\_date, is\_admin\}; \\
	email \rightarrow \{username, name, password, register\_date, is\_admin\}; \\
	name \rightarrow \{username, email, password, register\_date, is\_admin\}.
\end{array} \right.$ \\
\textit{Вывод:} все детерминанты ФЗ являются потенциальными ключами. Отношение в 3НФ.

\subsection*{Отношение games}

\begin{table}[H]
	\centering
	\caption{Атрибуты отношения games}
	\label{tab:games}
	\begin{tabularx}{\textwidth}{|l|X|X|}
		\hline
		\textbf{Атрибут} & \textbf{Тип} & \textbf{Назначение} \\
		\hline
		id (PK) & Строка & Идентификатор игры \\
		\hline
		name & Строка & Название игры \\
		\hline
		description & Текст & Описание проекта \\
		\hline
		image\_path & Строка & Путь к обложке \\
		\hline
	\end{tabularx}
\end{table}

\textit{Потенциальный ключ:} $K = \{id\}$. \\
\textit{ФЗ:} $F = \left\{ id \rightarrow \{name, description, image\_path\} \right.$ \\
\textit{Вывод:} неключевые атрибуты зависят только от ключа. Отношение в 3НФ.

\subsection*{Отношение mods}

\begin{table}[H]
	\centering
	\caption{Атрибуты отношения mods}
	\label{tab:mods}
	\begin{tabularx}{\textwidth}{|l|l|X|}
		\hline
		\textbf{Атрибут} & \textbf{Тип} & \textbf{Назначение} \\
		\hline
		id (PK) & Строка & Идентификатор модификации \\
		\hline
		name & Строка & Название модификации \\
		\hline
		description & Текст & Описание \\
		\hline
		image\_path & Строка & Путь к обложке \\
		\hline
		creation\_date & Дата и время & Дата добавления \\
		\hline
		game\_id (FK) & Строка & Ссылка на игру \\
		\hline
		publisher\_id (FK)& Строка & Ссылка на автора \\
		\hline
	\end{tabularx}
\end{table}

\textit{Потенциальный ключ:} $K = \{id\}$. \\
\textit{ФЗ:} $F = \left\{ id \rightarrow \begin{array}{l}
	\{name, description, image\_path, creation\_date, game\_id, \\
	publisher\_id\}
\end{array} \right.$ \\
\textit{Вывод:} транзитивные зависимости не обнаружены. Отношение в 3НФ.

\subsection*{Отношение mod\_versions}

\begin{table}[H]
	\centering
	\caption{Атрибуты отношения mod\_versions}
	\label{tab:mod_versions}
	\begin{tabularx}{\textwidth}{|l|X|X|}
		\hline
		\textbf{Атрибут} & \textbf{Тип} & \textbf{Назначение} \\
		\hline
		id (PK) & Универсальный уникальный идентификатор & Идентификатор версии \\
		\hline
		name & Строка & Название версии \\
		\hline
		changes & Текст & Список изменений \\
		\hline
		upload\_date & Дата и время & Дата загрузки \\
		\hline
		status & Перечисление & Статус проверки \\
		\hline
		mod\_id (FK) & Строка & Ссылка на мод \\
		\hline
	\end{tabularx}
\end{table}

\textit{Потенциальные ключи:} $K_1 = \{id\}$, $K_2 = \{name, mod\_id\}$. \\
\textit{ФЗ:} $F = \left\{ \begin{array}{l}
	id \rightarrow \{name, changes, upload\_date, status, mod\_id\}; \\
	\{name, mod\_id\} \rightarrow \{id, changes, upload\_date, status\}.
\end{array} \right.$ \\
\textit{Вывод:} детерминанты являются ключами. Отношение в 3НФ.

\subsection*{Отношение mod\_files}

\begin{table}[H]
	\centering
	\caption{Атрибуты отношения mod\_files}
	\label{tab:mod_files}
	\begin{tabularx}{\textwidth}{|l|X|l|}
		\hline
		\textbf{Атрибут} & \textbf{Тип} & \textbf{Назначение} \\
		\hline
		id (PK) & Универсальный уникальный идентификатор &  Идентификатор файла \\
		\hline
		filename & Строка & Имя файла \\
		\hline
		file\_path & Строка & Путь к файлу \\
		\hline
		downloads & Целое число число & Количество скачиваний \\
		\hline
		mod\_version\_id (FK)& Универсальный уникальный идентификатор & Ссылка на версию \\
		\hline
	\end{tabularx}
\end{table}

\textit{Потенциальный ключ:} $K = \{id\}$. \\
\textit{ФЗ:} $F = \left\{ id \rightarrow \{filename, file\_path, downloads, mod\_version\_id\}. \right.$ \\
\textit{Вывод:} поля зависят только от первичного ключа. Отношение в 3НФ.

\subsection*{Отношение categories}

\begin{table}[H]
	\centering
	\caption{Атрибуты отношения categories}
	\label{tab:categories}
	\begin{tabularx}{\textwidth}{|X|X|X|}
		\hline
		\textbf{Атрибут} & \textbf{Тип} & \textbf{Назначение} \\
		\hline
		name (PK) & Строка & Название категории \\
		\hline
		description & Текст & Описание категории \\
		\hline
	\end{tabularx}
\end{table}

\textit{Потенциальный ключ:} $K = \{name\}$. \\
\textit{ФЗ:} $F = \left\{ name \rightarrow \{description\}. \right.$ \\
\textit{Вывод:} отношение находится в 3НФ.

\subsection*{Отношение mods\_categories}

\begin{table}[H]
	\centering
	\caption{Атрибуты отношения mods\_categories}
	\label{tab:mods_categories}
	\begin{tabularx}{\textwidth}{|l|X|X|}
		\hline
		\textbf{Атрибут} & \textbf{Тип} & \textbf{Назначение} \\
		\hline
		mod\_id (PK, FK) & Строка & Идентификатор мода \\
		\hline
		category\_name (PK, FK)& Строка & Название категории \\
		\hline
	\end{tabularx}
\end{table}

\textit{Потенциальный ключ:} $K = \{mod\_id, category\_name\}$. \\
\textit{ФЗ:} $F = \emptyset$. \\
\textit{Вывод:} неключевые атрибуты отсутствуют. Отношение в 3НФ.

\subsection*{Отношение comments}

\begin{table}[H]
	\centering
	\caption{Атрибуты отношения comments}
	\label{tab:comments}
	\begin{tabularx}{\textwidth}{|l|X|X|}
		\hline
		\textbf{Атрибут} & \textbf{Тип} & \textbf{Назначение} \\
		\hline
		id (PK) & Универсальный уникальный идентификатор & Идентификатор комментария \\
		\hline
		content & Текст & Содержимое \\
		\hline
		publish\_date & Дата и время & Время отправки \\
		\hline
		user\_id (FK) & Строка & Автор \\
		\hline
		mod\_id (FK) & Строка & Модификация \\
		\hline
	\end{tabularx}
\end{table}

\textit{Потенциальный ключ:} $K = \{id\}$. \\
\textit{ФЗ:} $F = \left\{ id \rightarrow \{content, publish\_date, user\_id, mod\_id\} \right.$ \\
\textit{Вывод:} соблюдаются условия 3НФ.

\subsection*{Отношение collections}

\begin{table}[H]
	\centering
	\caption{Атрибуты отношения collections}
	\label{tab:collections}
	\begin{tabularx}{\textwidth}{|l|X|X|}
		\hline
		\textbf{Атрибут} & \textbf{Тип} & \textbf{Назначение} \\
		\hline
		id (PK) & Универсальный уникальный идентификатор & Идентификатор коллекции \\
		\hline
		name & Строка & Название \\
		\hline
		description & Текст & Описание подборки \\
		\hline
		owner (FK) & Строка & Владелец \\
		\hline
	\end{tabularx}
\end{table}

\textit{Потенциальные ключи:} $K_1 = \{id\}$, $K_2 = \{name, owner\}$. \\
\textit{ФЗ:} $F = \left\{ \begin{array}{l}
	id \rightarrow \{name, description, owner\}; \\
	\{name, owner\} \rightarrow \{id, description\}.
\end{array} \right.$ \\
\textit{Вывод:} левые части всех ФЗ являются ключами. Отношение в 3НФ.

\subsection*{Отношение collections\_mods}

\begin{table}[H]
	\centering
	\caption{Атрибуты отношения collections\_mods}
	\label{tab:collections_mods}
	\begin{tabularx}{\textwidth}{|l|X|l|}
		\hline
		\textbf{Атрибут} & \textbf{Тип} & \textbf{Назначение} \\
		\hline
		collection\_id (PK, FK)& Универсальный уникальный идентификатор & Идентификатор коллекции \\
		\hline
		order\_index (PK) & Целое число число & Позиция в списке \\
		\hline
		mod\_id (FK) & Строка & Идентификатор мода \\
		\hline
	\end{tabularx}
\end{table}

\textit{Потенциальные ключи:} $K_1 = \{collection\_id, order\_index\}$, \\$K_2 = \{collection\_id, mod\_id\}$. \\
\textit{ФЗ:} $F = \left\{ \begin{array}{l}
	\{collection\_id, index\} \rightarrow \{mod\_id\}; \\
	\{collection\_id, mod\_id\} \rightarrow \{order\_index\}.
\end{array} \right.$ \\
\textit{Вывод:} отношение находится в 3НФ.

\section{Разработка алгоритма работы триггера}

В системе пользователи загружают моды, которые состоят из версий и файлов. Версия мода может быть одобрена только в том случае, если к ней прикреплен хотя бы один файл.

Чтобы не допустить ситуацию, когда версия получает статус <<одобрено>>, но файлов у нее нет, был создан триггер. Он проверяет наличие файлов при изменении статуса версии и не дает сохранить некорректные данные.

Схема алгоритма работы триггера представлена на рисунке~\ref{img:trigger}.

\begin{figure}[H]
	\centering
	\vspace{1em}
	\includegraphics[width=0.7\textwidth]{trigger}
	\caption{Схема алгоритма работы триггера}
	\label{img:trigger}
\end{figure}

\section{Вывод}

В ходе проектирования базы данных была разработана схема, состоящая из 10 взаимосвязанных таблиц. С применением матричного подхода спроектировано распределение прав доступа между тремя ролями системы. Проведенный анализ функциональных зависимостей доказал, что все отношения строго удовлетворяют требованиям третьей нормальной формы, что гарантирует целостность данных и защиту от аномалий модификации. Был спроектирован алгоритм работы триггера для реализации ограничения логики модерации.